\textbf{I. Algebraic Solution}

% FIGURE NEEDED: sketch of the line of sight y = mx from the origin crossing
% the circular asteroid; 2008 theory solutions, p. 13, problem 13.
Let $y = mx$ be the line of sight of the telescope. \hfill(20\%)

Then, the intersection of the line of sight and the circle is
\begin{align*}
x^2 + (mx)^2 - 10x - 8(mx) + 40 &= 0\\
\left(1+m^2\right)x^2 - (10+8m)x + 40 &= 0
\end{align*}
\hfill(20\%)

The equation will have solution if its discriminant
\begin{align*}
D &= b^2 - 4ac \ge 0,\\
  &= (-(10+8m))^2 - 4(1+m^2)40 \ge 0.
\end{align*}
The solution of the inequality is
\[
\frac{5}{6} - \frac{\sqrt{10}}{12} \le m \le
\frac{5}{6} + \frac{\sqrt{10}}{12}
\]
\hfill(20\%)

Therefore the exact value of the maximum value of $\tan\varphi$ is
$\dfrac{5}{6} + \dfrac{\sqrt{10}}{12}$ \hfill(20\%)

and the minimum value is $\dfrac{5}{6} - \dfrac{\sqrt{10}}{12}$ \hfill(20\%)

\textbf{II. Geometric Solution}

% FIGURE NEEDED: circle of radius 1 centred at (5,4) with the two tangent
% lines from the origin, the line y = mx to the centre (angle alpha to the
% x-axis) and the half-opening angle beta between the centre line and a
% tangent; 2008 theory solutions, p. 13, problem 13.
Let us write the equation of the circle in the following form \hfill(20\%)
\begin{align*}
(x-5)^2 + (y-4)^2 &= -40 + 25 + 16\\
                  &= 1
\end{align*}
The centre of the circle is $(5,4)$ and the radius is 1.

First we obtain $\tan\alpha = \dfrac{4}{5}$ and
$\tan\beta = \dfrac{1}{\sqrt{40}}$ \hfill(20\%)

The minimum value of tan of elevation angle is \hfill(30\%)
\[
\tan(\alpha-\beta) = \frac{\tan\alpha - \tan\beta}{1 + \tan\alpha\tan\beta}
 = \frac{\dfrac{4}{5} - \dfrac{1}{\sqrt{40}}}{1 + \dfrac{4}{5\sqrt{40}}}
 = \frac{5}{6} - \frac{\sqrt{10}}{12}
\]
The maximum value of tan of elevation angle is
\[
\tan(\alpha+\beta) = \frac{\tan\alpha + \tan\beta}{1 - \tan\alpha\tan\beta}
 = \frac{\dfrac{4}{5} + \dfrac{1}{\sqrt{40}}}{1 - \dfrac{4}{5\sqrt{40}}}
 = \frac{5}{6} + \frac{\sqrt{10}}{12}
\]
\hfill(30\%)
